\section{From equality to pointwise matrix relations}
\label{comp:section}

We now apply Theorem~\ref{act:quantitative-action} to finite collections of
negative eigenfunctions.  Equality forces each of its nonnegative errors to
vanish.  The main issue is that the number of eigenfunctions may grow without
bound, whereas the potential has only $m$ channels.  We will use this fixed
number of columns to control the tails of the resulting matrices.

Throughout this section, $W$ is real symmetric, $W\geq0$ almost everywhere,
and $0<\int_{\R}\tr(W^p)<\infty$.  Suppose that equality holds, in the form
\begin{equation}
 b_\sigma\sum_{i\in\mathcal I}k_i^{2\gamma}
   =D_\sigma\int_{\R}\tr(W^p)\,dx,
 \label{comp:equality}
\end{equation}
where $-k_i^2$ are the negative eigenvalues of $H_W$, repeated according to
multiplicity, and
\[
 k_1\geq k_2\geq\cdots>0.
\]
By Lemma~\ref{pre:spectral}, the index set $\mathcal I$ is a nonempty finite
initial segment of $\mathbb N$ or all of $\mathbb N$, each eigenvalue has
finite multiplicity, and $k_i\to0$ in the infinite case.  Fix real
orthonormal eigenfunctions $\varphi_i\in H^1(\R;\R^m)$ and regard them as
row vectors.  Write
\[
 \mathscr H=\ell^2(\mathcal I;\R),\qquad
 \Lambda=\operatorname{diag}(k_i:i\in\mathcal I),\qquad \kappa=k_1.
\]
For operators between real Hilbert spaces, the superscript $T$ will also
denote the adjoint.

\begin{proposition}[Pointwise relations at equality]
\label{comp:relations}
Under these assumptions, there are rows $a_i\in H^1(\R;\R^m)$,
$i\in\mathcal I$, each a linear combination of eigenfunctions with
eigenvalue $-k_i^2$, such that
\begin{equation}
 -a_i''-a_iW=-k_i^2a_i
 \quad\text{in distributions.}
 \label{comp:row-eigenfunction}
\end{equation}
The rows $a_i$ are allowed to vanish.  On a common set of full measure, they
define a Hilbert--Schmidt operator $A(x):\R^m\to\mathscr H$, and there exists
a bounded self-adjoint operator $B$ on $\mathscr H$ with
$\|B\|_{\op}\leq\kappa$ such that
\begin{equation}
 [B,\Lambda^2]=0,\qquad
 \Lambda^2-B^2=(AA^T)^r,\qquad
 a_i'(x)=(BA)_{i,\cdot}\quad(i\in\mathcal I).
 \label{comp:limit-relations}
\end{equation}
Moreover,
\begin{equation}
 W(x)=q\bigl(A(x)^TA(x)\bigr)^r
 \quad\text{almost everywhere.}
 \label{comp:potential-limit}
\end{equation}
The operator $B$ is asserted to exist separately at each such point $x$;
no measurable choice is required.
\end{proposition}

We divide the proof into the construction of a sequence whose errors
vanish, convergence of every fixed row, and control of the remaining rows.

\subsection{Finite spectral approximations and vanishing errors}
\label{comp:approximation}

Let $n_j=j$ when $\mathcal I=\mathbb N$, and let $n_j=|\mathcal I|$ when
$\mathcal I$ is finite.  Choose $0<\xi_j<1/j$ so small that
\begin{equation}
 \ell_{j,i}:=(k_i^2-\xi_j)^{1/2}>0\quad(1\leq i\leq n_j),
 \qquad
 \sum_{i\leq n_j}\bigl(k_i^{2\gamma}-\ell_{j,i}^{2\gamma}\bigr)<1/j.
 \label{comp:common-shift}
\end{equation}
Set $K_j=\operatorname{diag}(\ell_{j,1},\ldots,\ell_{j,n_j})$.
The same shift $\xi_j$ is used for every eigenvalue: this leaves the
differences of their squares unchanged.

The approximation leading to \eqref{bound:trial} can be made arbitrarily
accurate in $H^1$.  Apply it to the first $n_j$ eigenfunctions, choosing
the $H^1$ errors below $1/j$ and the form-matrix error below $\xi_j$.
This gives real, smooth, compactly supported functions
$u_{j,1},\ldots,u_{j,n_j}$, orthonormal in $L^2(\R;\R^m)$, such that
\begin{equation}
 \|u_{j,i}-\varphi_i\|_{H^1}<1/j,
 \qquad
 \bigl(h_W[u_{j,i},u_{j,l}]\bigr)_{i,l=1}^{n_j}\leq-K_j^2.
 \label{comp:form-approximation}
\end{equation}
Here $h_W[\cdot,\cdot]$ is the polarized quadratic form, and the second
inequality is the ordering of real symmetric matrices.

Let $U_j$ be the matrix with rows $u_{j,i}$, and choose a bounded interval
$\Omega_j$ containing their supports in its interior.  Apply
Theorem~\ref{act:quantitative-action} with $K=K_j$, $d=m$, $\kappa=k_1$,
and $\eta=1/j$.  It gives a lower triangular matrix $C_j$, a symmetric
path $B_j$ with $\|B_j\|_{\op}\leq\kappa$, and $A_j=C_jU_j$.
Extend $A_j$ by zero outside $\Omega_j$ and extend $B_j$ by $K_j$ to the
left and by $-K_j$ to the right.  All the errors in that theorem vanish
outside $\Omega_j$, so their integrals may be taken over $\R$.

The lower triangular form of $C_j$ gives a useful additional error.  Put
$Z_j=A_j^TA_j$, an $m\times m$ positive semidefinite matrix.  By
\eqref{comp:form-approximation},
\begin{equation}
 \mathcal E_{K_j}(A_j)
 \leq\int_{\R}\bigl(\tr(WZ_j)-\tr(Z_j^q)\bigr)\,dx-T_j,
 \qquad
 T_j=\sum_{l\leq i\leq n_j}(C_j)_{il}^{\,2}(k_l^2-k_i^2)\geq0.
 \label{comp:triangular-gap}
\end{equation}
Indeed, the $i$th row of $A_j$ is $\sum_{l\leq i}(C_j)_{il}u_{j,l}$.
Its quadratic form, plus $\ell_{j,i}^2$ times its squared $L^2$ norm, is
at most
\[
 \sum_{l\leq i}(C_j)_{il}^{\,2}
      (\ell_{j,i}^2-\ell_{j,l}^2)
 =-\sum_{l\leq i}(C_j)_{il}^{\,2}(k_l^2-k_i^2).
\]
Summing over $i$ proves \eqref{comp:triangular-gap}, since $A_jA_j^T$ and
$Z_j$ have the same nonzero eigenvalues, with multiplicity.  The ordering
of the $k_i$ makes every summand of $T_j$ nonnegative.

Recall the nonnegative Young deficit from Lemma~\ref{pre:young}:
\begin{equation}
 Y(W,Z)=D_\sigma\tr(W^p)-\tr(WZ)+\tr(Z^q)\geq0,
 \qquad
 Y(W,Z)=0\ \Longleftrightarrow\ W=qZ^r.
 \label{comp:young-deficit}
\end{equation}
Combining Theorem~\ref{act:quantitative-action},
\eqref{comp:equality}, and \eqref{comp:triangular-gap} yields
\begin{align}
 0\leq T_j+\int_{\R}\Bigl(
   Y(W,Z_j)+\|A_j'-B_jA_j\|_{\HS}^2
       +c_*\Phi_{K_j}(A_j,B_j)\Bigr)\,dx
 &\leq d_j,\label{comp:vanishing-errors}\\
 d_j:=b_\sigma\left(
    \sum_{i\in\mathcal I}k_i^{2\gamma}
       -\sum_{i\leq n_j}\ell_{j,i}^{2\gamma}\right)+1/j
 &\longrightarrow0.\notag
\end{align}
Here $c_*>0$ is fixed: its independence of $n_j$ is essential.  Thus all
three integral errors tend to zero, as does $T_j$.

We will also need a uniform bound on $A_j$.  The coercive estimate
\eqref{pre:young-coercive} gives
\begin{equation}
 Y(W,Z_j)\geq\frac12\tr(Z_j^q)
       -(2^{p-1}-1)D_\sigma\tr(W^p).
 \label{comp:coercivity}
\end{equation}
Since $Z_j$ has size $m$, scalar H\"older's inequality gives
\[
 \|A_j\|_{\HS}^{2q}=(\tr Z_j)^q
       \leq m^{q-1}\tr(Z_j^q).
\]
Equations~\eqref{comp:vanishing-errors} and \eqref{comp:coercivity}
therefore bound $A_j$ uniformly in $L^{2q}(\R)$ with the
Hilbert--Schmidt norm, independently of its number of rows.

\subsection{Convergence of each fixed row}
\label{comp:rows}

For fixed $i\in\mathcal I$, the coefficients $(C_j)_{il}$, $l\leq i$,
are bounded as $j\to\infty$.  To see this, choose a bounded interval $J_i$
so large that the restricted Gram matrix
\[
 \left(\int_{J_i}\varphi_l(x)\varphi_s(x)^T\,dx
       \right)_{l,s=1}^{i}\geq\frac34 I_i,
\]
where $I_i$ denotes the $i\times i$ identity matrix.
Such an interval exists because the full Gram matrix is the identity.
By \eqref{comp:form-approximation}, the corresponding restricted Gram
matrix of the $u_{j,l}$ is at least one half of the identity for all
sufficiently large $j$.  Hence
\[
 \frac12\sum_{l\leq i}(C_j)_{il}^{\,2}
 \leq\int_{J_i}|(A_j)_{i,\cdot}|^2\,dx
 \leq |J_i|^{1-1/q}
        \left(\int_{J_i}\|A_j\|_{\HS}^{2q}\,dx\right)^{1/q}.
\]
The last expression is uniformly bounded.  Lower triangularity is used
here to keep the number of coefficients in a fixed row finite.

A diagonal subsequence now makes every coefficient $(C_j)_{il}$ converge
to a limit $c_{il}$, for $l\leq i$.  From
$T_j\to0$ and \eqref{comp:triangular-gap},
\begin{equation}
 c_{il}=0\quad\text{whenever }k_l>k_i.
 \label{comp:gap-suppression}
\end{equation}
Define
\[
 a_i=\sum_{l\leq i}c_{il}\varphi_l.
\]
The sum is finite, and \eqref{comp:form-approximation} implies
\begin{equation}
 (A_j)_{i,\cdot}\longrightarrow a_i
       \quad\text{in }H^1(\R;\R^m)
       \quad\text{for every fixed }i.
 \label{comp:row-convergence}
\end{equation}
The ordering of the eigenvalues and \eqref{comp:gap-suppression} show
that only eigenfunctions with eigenvalue $-k_i^2$ occur in $a_i$.
This proves \eqref{comp:row-eigenfunction}, including the possibility
$a_i=0$.  The eigenfunction equations hold distributionally by testing
the form equations against smooth compactly supported functions; the
potential is locally integrable and the rows are locally bounded.

Choose the continuous, locally absolutely continuous representatives
of these rows.  The Sobolev embedding $H^1(\R)\subset C_b(\R)$ shows
that \eqref{comp:row-convergence} also gives pointwise convergence at
every $x$.  Passing to a further subsequence, we may arrange that the
integrals of the nonnegative integrand in
\eqref{comp:vanishing-errors} are summable, and that, for every fixed
$i$, the squared $L^2$ errors of the derivatives in
\eqref{comp:row-convergence} are summable.  A diagonal choice achieves
both conditions, by making the first $j$ available derivative errors and the
integral error less than $2^{-j}$ along the chosen subsequence.
Consequently, on a common set of full measure,
\begin{equation}
 \begin{gathered}
 Y(W,Z_j)\to0,\qquad
 \|A_j'-B_jA_j\|_{\HS}\to0,\qquad
 \Phi_{K_j}(A_j,B_j)\to0,\\
 (A_j')_{i,\cdot}(x)\to a_i'(x)
       \quad\text{for every }i\in\mathcal I.
 \end{gathered}
 \label{comp:pointwise-errors}
\end{equation}
We also exclude the null set where $W$ is not a finite positive
semidefinite matrix.  It remains to show that convergence of the
individual rows is strong enough to pass to the products in these
relations.

\subsection{Tail control in the row-index space}
\label{comp:tails}

Pad $A_j$, $B_j$, and $K_j$ with zeros in the unused coordinates of
$\mathscr H$.  Thus $A_j:\R^m\to\mathscr H$, and $B_j,K_j$ are
self-adjoint finite-rank operators on $\mathscr H$.  Since the omitted
$k_i$ tend to zero,
\begin{equation}
 \|K_j^2-\Lambda^2\|_{\op}
    \leq\max\{\xi_j,\sup_{i>n_j}k_i^2\}\longrightarrow0,
 \label{comp:diagonal-limit}
\end{equation}
where a supremum over no indices is zero.

Fix a point $x$ in the full-measure set from
\eqref{comp:pointwise-errors} and suppress $x$ for the moment.  Write
$\rho_j=A_jA_j^T$.  The definition of $\Phi$ and its convergence to zero
give
\begin{equation}
 \|[B_j,K_j^2]\|_{\op}\to0,
 \qquad e_j:=\|K_j^2-B_j^2-\rho_j^r\|_{\op}\to0.
 \label{comp:pointwise-algebra-error}
\end{equation}
For $h\geq1$, let $P_h$ be the orthogonal projection onto the first $h$
coordinates of $\mathscr H$, with all coordinates included if $h$ exceeds
$|\mathcal I|$.  Set $Q_h=I-P_h$ and
$s_h=\sup_{i>h}k_i^2$, again with the empty supremum equal to zero.
Then $s_h\to0$.  For every unit vector $v$ in the range of $Q_h$,
\begin{equation}
 \|B_jv\|^2+\langle v,\rho_j^rv\rangle\leq s_h+e_j.
 \label{comp:tail-quadratic}
\end{equation}
Both terms on the left are nonnegative.  Jensen's inequality for the
spectral measure of $\rho_j$, using $r>1$, implies
\[
 \langle v,\rho_jv\rangle^r\leq\langle v,\rho_j^rv\rangle.
\]
It follows that
$\|Q_h\rho_jQ_h\|_{\op}\leq(s_h+e_j)^{1/r}$.
Although $\mathscr H$ may be infinite dimensional, $Q_h\rho_jQ_h$
has rank at most $m$, because $A_j$ has exactly $m$ columns.  Its trace
is therefore at most $m$ times its operator norm.  We obtain the two
tail bounds
\begin{equation}
 \boxed{\quad
 \|Q_hA_j\|_{\HS}^2\leq m(s_h+e_j)^{1/r},
 \qquad
 \|B_jQ_h\|_{\op}\leq(s_h+e_j)^{1/2}.
 \quad}
 \label{comp:tail-bounds}
\end{equation}
The first estimate converts decay of the eigenvalue parameters into
decay of the total mass of all omitted rows.  Its factor is the fixed
physical dimension $m$, rather than the number of retained
eigenfunctions; Figure~\ref{comp:tail-figure} depicts this distinction.

\begin{figure}[tb]
 \centering
 \begin{tikzpicture}[font=\small,>=Latex]
  \fill[black!5] (0,2.05) rectangle (2.7,3.65);
  \fill[blue!7] (0,0) rectangle (2.7,2.05);
  \draw[thick] (0,0) rectangle (2.7,3.65);
  \draw[thick] (0,2.05)--(2.7,2.05);
  \draw[<->] (0,4.05)--node[above]{$m$ columns, fixed}(2.7,4.05);
  \node at (1.35,2.83) {$P_hA_j$};
  \node at (1.35,1.28) {$Q_hA_j$};
  \node at (1.35,0.55) {$\vdots$};
  \node[anchor=east,align=right] at (-0.2,2.83)
       {first $h$\\rows};
  \node[anchor=east,align=right] at (-0.2,1.08)
       {all remaining\\rows};
  \draw[->] (2.9,2.83)--(3.6,2.83);
  \node[anchor=west,align=left] at (3.75,2.83)
       {Finite-dimensional convergence\\for each fixed $h$};
  \draw[->] (2.9,1.08)--(3.6,1.08);
  \node[anchor=west,align=left] at (3.75,1.08)
       {$\operatorname{rank}(Q_hA_j)\leq m$\\[4pt]
        $\|Q_hA_j\|_{\HS}^2\leq m(s_h+e_j)^{1/r}$};
 \end{tikzpicture}
 \caption{The number of rows can grow with the number of negative
 eigenfunctions, but the number of physical channels remains $m$.
 After fixing the first $h$ rows, the identity
 $K_j^2\approx B_j^2+(A_jA_j^T)^r$ bounds the operator norm in the
 remaining rows.  The rank bound converts it into a Hilbert--Schmidt
 bound, which tends to zero as first $j\to\infty$ and then $h\to\infty$.
 The drawing is schematic; all matrices are padded by zero rows in
 $\mathscr H$.}
 \label{comp:tail-figure}
\end{figure}

For each fixed $h$, the first $h$ rows of $A_j$ converge by
\eqref{comp:row-convergence}.  The first bound in
\eqref{comp:tail-bounds} makes the remaining rows uniformly small by
choosing $h$ large and then $j$ sufficiently large.  Thus $A_j$ is Cauchy in
Hilbert--Schmidt norm and converges to an operator $A:\R^m\to\mathscr H$
whose rows are precisely $a_i(x)$:
\begin{equation}
 A_j\longrightarrow A\quad\text{in Hilbert--Schmidt norm.}
 \label{comp:HS-limit}
\end{equation}
This argument applies to the whole subsequence already selected and
does not require a further choice depending on $x$.

For $B_j$ we need only pointwise compactness.  Its finite compressions
$P_hB_jP_h$ are bounded in finite-dimensional spaces, so a diagonal
subsequence makes them converge for every $h$.  Self-adjointness and
the second estimate in \eqref{comp:tail-bounds} give
\[
 \|B_j-P_hB_jP_h\|_{\op}
   \leq2\|B_jQ_h\|_{\op}
   \leq2(s_h+e_j)^{1/2}.
\]
The chosen subsequence consequently converges in operator norm to a
self-adjoint $B$ with $\|B\|_{\op}\leq\kappa$.  This subsequence, and
hence $B$, may depend on the point $x$.

We can now pass to all the required products.  Hilbert--Schmidt
convergence of $A_j$ implies operator norm convergence of
$A_jA_j^T$ and of $A_j^TA_j$.  Continuous functional calculus for
nonnegative operators then gives convergence of their $r$th powers.
Together with \eqref{comp:diagonal-limit} and
\eqref{comp:pointwise-algebra-error}, this proves
\[
 [B,\Lambda^2]=0,\qquad
 \Lambda^2-B^2=(AA^T)^r.
\]
Also $B_jA_j\to BA$ in Hilbert--Schmidt norm.  Taking its $i$th row
and using both the derivative convergence and the first-order error
in \eqref{comp:pointwise-errors} gives
$a_i'(x)=(BA)_{i,\cdot}$ for every $i$.  Finally, continuity of the
finite-dimensional Young deficit and
$Y(W,A_j^TA_j)\to0$ yield
$Y(W,A^TA)=0$.  Lemma~\ref{pre:young} gives
$W=q(A^TA)^r$.

We have proved these statements at every point of one common set of
full measure.  The construction selects no function $x\mapsto B(x)$:
only the existence of a suitable operator at each such point enters
the next section.  This completes the proof of
Proposition~\ref{comp:relations}.
